\section{Geometric estimates for finite two-command acquisition}
\label{sec:two-field-geometric-details}
This section gives the uniform geometric estimates used in
Section~\ref{sec:two-field-finite}.  The constants depend only on
\eqref{eq:two-field-geometric-priors},
\eqref{eq:two-field-finite-priors}, and the fixed positive length
$t$.  In particular they do not depend on a modulus of continuity
of the unknown density.  Put
\[
 r_- =\kappa_+^{-1},\qquad r_+=\kappa_-^{-1},\qquad
 g=d_0-\Delta>2t,\qquad D_P=D_0+\Delta.
\]
An expanded component has curvature radius between
$r_-+r_A^-$ and $r_++r_A^+$, diameter at most $D_P$, and
distance at least $g$ from every other expanded component.

\subsection{A fixed collar and the conjunction of two tests}
\begin{lemma}[Quantitative tangential collar]
\label{lem:two-field-collar-details}
Choose $R>\max(t,r_++r_A^+)$, and fix
\[
 0<d_c<\min\left\{\frac{r_-+r_A^-}{4},\frac R2,
                 \frac{t^2}{32R},g-t\right\},
 \qquad 0<\eta_n\le\frac{t}{48R}.
\]
Let $y=y_0-dn$, where $y_0\in\partial P_C$, $n$ is its outward
unit normal, and $|d|\le d_c$.  Suppose that a rational vector
$\widehat n$ satisfies $|\widehat n-n|\le\eta_n$.  Define
\[
 V=\left\{v\in\{e_1,-e_1\}:\widehat n\cdot v
     \ge\max_{w\in\{e_1,-e_1\}}\widehat n\cdot w-4\eta_n\right\}.
\]
For every $v\in V$, use nominal center $y+tv$ and displacement
$-tv$.  If $d<0$, at least one of these experiments has collision
probability zero.  If $e\le d\le d_c$, with $0<e<e_0$ for a fixed
$e_0>0$, every one has collision probability at least
$q_e=c e^{\gamma+3/2}$, where $c>0$ is known from the priors.
The rule which returns one if and only if each of the at most two
batches contains a collision therefore has error probability at most
$\varepsilon$ outside the inner layer of width $e$, using at most
\[
 2\left\lceil q_e^{-1}\log(2/\varepsilon)\right\rceil
\]
attempts.  These statements hold conditionally on successful coarse
normal acquisition and on all subsequent adaptive history.
\end{lemma}
\begin{proof}
The lower rolling radius makes the signed normal coordinates unique
in the displayed collar.  The upper curvature-radius bound gives
\[
             P_C\subset\overline B_R(y_0-Rn).
\]
For clarity, if the normal angle of $n$ is zero, the support of
$P_C-y_0$ equals
$\int_0^\phi r_{P_C}(u)\sin(\phi-u)\dd u$ for
$0\le\phi\le\pi$, and is at most $R(1-\cos\phi)$.
Integrating along the reverse angular interval gives the same bound
for $-\pi\le\phi\le0$.  These are the support inequalities for
the asserted disk.

Every true maximizer $v_*$ of $n\cdot v$ over the two directions
belongs to $V$: replacing $n$ by $\widehat n$ changes a comparison
by at most $2\eta_n$.  For every $v\in V$,
\[
             n\cdot v\ge n\cdot v_*-6\eta_n
                           \ge-\frac{t}{8R}.
\]
Consequently, if $0<d\le d_c$,
\begin{align*}
 |y+tv-(y_0-Rn)|^2-R^2
  &=t^2+2t(R-d)n\cdot v-2Rd+d^2\\
  &\ge t^2-\frac{t^2}{4}-\frac{t^2}{16}
      =\frac{11t^2}{16}>0.
\end{align*}
Thus $y+tv\notin P_C$, including when the command is nearly tangent
to $\partial P_C$.  Since $P_C=C-A$, every physical start
$y+tv+z$, $z\in A$, lies outside the closed obstacle $C$.

Write $y_0=c_0-z_0$ for the support points of $C$ and $A$ with
normals $n$ and $-n$.  Every point of every candidate segment has
the form
\[
          c_0+(z-z_0)-dn+(1-u)tv,
                   \qquad z\in A,\quad 0\le u\le1.
\]
Its distance from $c_0$ is at most $\Delta+d_c+t<d_0$.
Thus no other obstacle meets any candidate segment.  When $d<0$,
the true maximizing candidate has $n\cdot v_*\ge0$, and
\[
 n\cdot\{y+tv_*+z-utv_*\}
 =h_C(n)-d+n\cdot(z-z_0)+(1-u)t n\cdot v_*
 >h_C(n).
\]
It always misses.  This proves the exterior assertion for the
conjunction; individual nonmaximizing candidates need not have zero
collision probability there.

For positive depth, the event $y+Z\in C$ forces a collision from
every free shifted start just established.  The overlap estimate
of Lemma~\ref{H-lem:cap-mass} gives probability at least
$q_e=c e^{\gamma+3/2}$.  Its constant is uniform here: the proof
uses only lower rolling radii, upper diameter and support bounds,
and $b_0,\gamma$.  The footprint $K$ in that lemma may therefore
range over the present class of unknown bodies $A$.  More
explicitly, replacing $A$ by $A_{-e/4}$ leaves a uniform rolling
radius, gives an overlap of area at least $c e^{3/2}$, and on that
overlap the density is at least $b_0(e/4)^\gamma$.

Fresh preparations have the same law at every candidate.  A batch
of the stated size has no collision with probability at most
$(1-q_e)^{\lceil q_e^{-1}\log(2/\varepsilon)\rceil}
\le\varepsilon/2$.  The union bound over $V$ proves the assertion.
The coarse event fixes the validity of the normal estimate; fresh
draws then prove the same bound after conditioning on the entire
previous adaptive record.  For a rounded nominal target the lemma
is applied at that actual target, which must stay in the fixed
normal neighborhood.  The signed depth is a Lipschitz function,
so rounding by $O(e)$ enlarges the intended radial bisection
ambiguity by $O(e)$.  The sure-zero exterior assertion refers to
the actual rounded target.
\end{proof}

\subsection{Area at grazing points of a collision strip}
\begin{lemma}[A uniform strip neighborhood]
\label{lem:two-field-strip-details}
For either sign let $S_C^\pm$ be the closed collision strip of
Section~\ref{sec:two-field-finite}.  There are known constants
$c_S,r_S>0$ such that
\[
        \area(S_C^\pm\cap B(y,r))\ge c_S r^3,
              \qquad y\in S_C^\pm,\quad 0<r<r_S.
\]
This estimate includes both tangent endpoints and both ends of each
flight segment.
\end{lemma}
\begin{proof}
It suffices to consider the positive direction.  In normal angle,
\[
 F(\phi,u)=c_C(\phi)-u e_1,
   \qquad \frac\pi2\le\phi\le\frac{3\pi}2,
   \quad 0\le u\le t.
\]
The transverse coordinate $c_C(\phi)\cdot e_2$ is strictly
decreasing on the open angular interval because its derivative is
$r_C(\phi)\cos\phi<0$.  The parametrization is consequently
injective there.  Its absolute Jacobian is
$r_C(\phi)|\cos\phi|$, and its image is the strip, apart from
boundary curves of area zero.

Fix $y=F(\phi_0,u_0)$, and reduce $r_S$ so that
$r<t$ and
\[
 \alpha=\frac{r}{16(1+r_+)}<\frac{\pi}{12},
       \qquad v_r=\frac r{16}<\frac t4.
\]
Move $\phi_0$, if necessary, to the closest point $\phi_*$ of
$[\pi/2+2\alpha,3\pi/2-2\alpha]$, and put
$I=[\phi_*-\alpha/2,\phi_*+\alpha/2]$.
Then $|I|=\alpha$, $|\phi-\phi_0|\le5\alpha/2$ on $I$,
and $|\cos\phi|\ge\alpha/2$ there.  Similarly move $u_0$
to the closest point $u_*$ of $[v_r,t-v_r]$ and put
$J=[u_*-v_r/2,u_*+v_r/2]$.
Thus $|J|=v_r$ and $|u-u_0|\le3v_r/2$.

Since $|c_C'(\phi)|\le r_+$, every point of $F(I\times J)$
is at distance at most
\[
        \frac52r_+\alpha+\frac32v_r<r
\]
from $y$.  The area formula on this injective rectangle now gives
\[
 \area(S_C^+\cap B(y,r))
 \ge\int_I\int_J r_C(\phi)|\cos\phi|\dd u\dd\phi
 \ge\frac{r_-}{2}\alpha^2v_r=c_Sr^3.
\]
The construction deliberately moves the angular rectangle into the
incoming half-circle.  At a tangent endpoint the Jacobian loses one
power of $r$, which is exactly the third power in the bound.
Reflection in the first coordinate proves the negative case.
\end{proof}

\subsection{Finite nominal grids and complete components}
\begin{lemma}[Boundary cells and positive cap probabilities]
\label{lem:two-field-grid-details}
Let $W$ be a fixed dyadic rectangle containing
$P_C+t\overline B$ for every component to be retained, with a
fixed positive collar.  Put
$\kappa=\gamma+9/2$ and
$Q_C^\pm=\operatorname{conv}(S_C^\pm)-A$.
For every unit direction $u$ and sufficiently small $\epsilon>0$,
a uniform nominal center in $W$, followed by the corresponding
fixed command, has probability at least $c\epsilon^\kappa$ of
producing a positive record from $C$ in the support cap
\[
 \{x\in Q_C^\pm:
       u\cdot x\ge h_{Q_C^\pm}(u)-\epsilon\}.
\]
The same assertion holds, with a smaller $c$, when the nominal center
is uniform among the centers of the square grid cells of side
$\ell$ partitioning $W$, provided
$\ell\le c_0\epsilon^\kappa$.  No bound on the density other
than \eqref{eq:two-field-finite-priors} is used.
\end{lemma}
\begin{proof}
Choose a point $y_u\in S_C^\pm$ maximizing $u\cdot y$.
Lemma~\ref{lem:two-field-strip-details}, applied with radius
$\epsilon/6$, gives strip area at least $c\epsilon^3$ with
support deficit at most $\epsilon/6$.

At the footprint support point $z_0$ with outward normal $-u$,
the interior disk of radius $r_A^-$ lies in $A$.  With $v=-u$,
use coordinates
\[
 z=z_0-qv+w v^\perp,\qquad
 \frac\epsilon8\le q\le\frac\epsilon6,
 \qquad |w|\le\frac{\sqrt{r_A^-\epsilon}}{16}.
\]
For $\epsilon$ sufficiently small relative to $r_A^-$ this
rectangle has distance at least $\epsilon/64$ from the boundary
of the interior disk, hence from $\partial A$.  For example its
squared distance from the disk center is
$(r_A^--q)^2+w^2\le(r_A^--\epsilon/64)^2$ after decreasing
the fixed upper bound for $\epsilon$.
Its area is $c\epsilon^{3/2}$ and its support deficit in
direction $-u$ is $q\le\epsilon/6$.  Its probability mass is
therefore at least $c\epsilon^{\gamma+3/2}$.

For every pair $(y,z)$ in these two subsets, $x=y-z$ belongs to
$W$ and has support deficit at most $\epsilon/3$ in $Q_C^\pm$.
Fubini's theorem, divided by $|W|$, proves the continuous lower
bound $c\epsilon^{\gamma+9/2}$.  Strip boundary points can be
removed here: they have zero area in the integration over $y$.

We give the grid argument with its conditioning order explicit.
Fix $Z=z$.  Membership of a nominal center in the relevant event
is determined by a translated collision strip, one half-plane,
and $W$.  The condition $x\in Q_C^\pm$ adds no boundary, since
$x\in S_C^\pm-z$ and $z\in A$ already imply it.
Up to boundary conventions the strip indicator is the difference
of the indicators of $C+[\mp te_1,0]$ and $C$.
Each of these convex sets has a uniformly bounded perimeter; for
example its perimeter is at most $\pi(D_0+t)$.
The additional half-plane cuts and the sides of $W$ have bounded
total length.

Only grid cells meeting one of these boundaries can disagree with
the value of the indicator at the cell center.  A curve of length
$L$ meets cells of total area at most
$C(L\ell+\ell^2)$: partition the curve by arc length into
at most $1+L/\ell$ subarcs of length at most $\ell$; each
subarc and all cells meeting it are contained in a disk of radius
$3\ell$ about its initial point.  This also covers cells whose
centers lie exactly on a boundary.  Thus the discrepancy between
the uniform continuous probability and the uniform grid probability
is at most $C\ell$, uniformly in the fixed $z$.
Integration against the probability law of $Z$ preserves this
bound.  Taking $\ell\le c_0\epsilon^\kappa$ with $c_0$
small enough preserves half the continuous lower bound.
\end{proof}

\begin{lemma}[Assignment and complete hull acquisition]
\label{lem:two-field-assignment-details}
Suppose that the successful coarse event supplies complete expanded
components $P_C$ with convex hull estimates $H_C$ satisfying
$d_H(H_C,P_C)\le\sigma$, where
\[
            0<\sigma\le\frac{g-2t}{16}.
\]
A positive nominal record is assigned to $C$ if the upper endpoint
of a distance enclosure for $\dist(x,H_C)$ is less than
$t+2\sigma$, using an enclosure of width at most $\sigma/4$.
This rule assigns all
records from $C$ correctly and never assigns another component's
record to $C$.  It remains valid when only the required complete
components are retained and all other records are discarded.

With the grid of Lemma~\ref{lem:two-field-grid-details}, the convex
hulls of the assigned positive nominal centers have simultaneous
Hausdorff error at most $C\epsilon$ after
\[
       C\epsilon^{-\kappa}\log\frac C{\epsilon\delta}
\]
attempts, for the two signs and all of the prior-bounded number of
required complete components, with conditional probability at least
$1-\delta$.
\end{lemma}
\begin{proof}
A positive record $x$ from $C$ belongs to $S_C^\pm-A$, so
$\dist(x,P_C)\le t$.  Hence
\[
 \dist(x,H_C)\le t+\sigma.
\]
If the same record belongs instead to a different component $D$,
then
\[
 \dist(x,H_C)\ge g-t-\sigma
       \ge t+15\sigma.
\]
The threshold has a reserve of $\sigma$ on the first side and
$13\sigma$ on the second.  The specified numerical enclosure
preserves both decisions.  A segment of length $t<d_0$ cannot
meet two distinct obstacles, so every positive record has a unique
physical component.  That component label is established by this
distance test; it is not a returned part of the observation.

Here is a concrete completeness cutoff.  Perform the coarse
occupation search on a larger square $U$, with grid spacing and
depth thresholds fixed from the priors.  On its accuracy event all
retained nodes lie in their assigned $P_C$, all sufficiently deep
nodes are retained, and the clustering and interior-ball argument
of Lemma~\ref{H-lem:coarse} applies.  Discard every cluster hull
whose distance from $\partial U$ is at most $D_P+2\sigma$.
If $P_C$ meets the exterior of $U$, select a point of $P_C$ on
one exterior side.  The diameter bound $D_P$ forces every retained
node of this component, and its entire cluster hull, to lie in the
width-$D_P$ strip along that side of $U$.  Such a fragment is
discarded.  Thus every surviving hull represents a complete
component.  Taking $U$ with a collar exceeding $2D_P+4\sigma$
about the required region protects its complete components from
this cutoff.  Finally enlarge $W$, still by a fixed amount, to
contain their $P_C+t\overline B$ and a positive collar.  All
their possible nominal collision records are then sampled in
$W$.  Fragments of other components cannot pass the distance
test just proved, whether or not those fragments have been
retained in the coarse list.

For each required $C$ and sign, take a direction net of spacing
$c\epsilon$ on the circle, of cardinality at most
$C\epsilon^{-1}$.  The net is used only in the proof.  The
probability that a particular cap receives no assigned positive
record in $N$ independent attempts is at most
$\exp(-cN\epsilon^\kappa)$ by
Lemma~\ref{lem:two-field-grid-details} and the correct assignment.
The union bound over the net, components, and signs gives a record
in every such cap at the asserted $N$.

Every retained center lies in $Q_C^\pm$, so its convex hull is
inscribed.  Both the true and empirical supports have a uniform
Lipschitz constant in the direction, since the corresponding sets
lie in the fixed bounded rectangle $W$.  Extending the cap
inequalities from the net therefore gives supremum support error
$C\epsilon$, equivalently Hausdorff error $C\epsilon$.
All retained vertices are known dyadic nominal centers.  Neither
collision locations nor a sample of the launch displacement is
needed for their rational polygon representation.
\end{proof}

\subsection{Signed plateaux and curvature-atom mass}
\begin{lemma}[A fixed plateau with zero moments]
\label{lem:two-field-plateau-details}
There is an even $C^\infty$ function $\psi$, supported in
$(-1,1)$ and equal to one on $[-1/8,1/8]$, such that
\[
             \int u^k\psi(u)\dd u=0,
                         \qquad 0\le k\le4.
\]
The function can be fixed once by an invertible three-dimensional
linear system.  If
\[
 (\partial_\phi^2+1)H
       =r(\phi)\dd\phi-\ell(\delta_\theta+\delta_{\theta+\pi}),
       \qquad \|r\|_{C^{4,\beta}}\le M,
\]
$\|\widehat H-H\|_\infty\le\epsilon$, and
$|\widehat\theta-\theta|\le C_0\epsilon$, then, with
$b\asymp\epsilon^{1/(6+\beta)}$,
\[
 -\int\widehat H(\phi)
   \left\{b^{-2}\psi''\left(\frac{\phi-\widehat\theta}{b}\right)
             +\psi\left(\frac{\phi-\widehat\theta}{b}\right)\right\}
             \dd\phi
       =\ell+O(\epsilon^{(5+\beta)/(6+\beta)}).
\]
The implied constant depends only on $M,C_0$ and the fixed function
$\psi$.
\end{lemma}
\begin{proof}
Take an even smooth central function $\psi_0$, equal to one on
$[-1/8,1/8]$ and supported in $(-1/4,1/4)$.  Fix an even
nonnegative smooth function $k$, supported in $(-1,1)$ and with
integral one.  Put
\[
 x_1=\frac13,\qquad x_2=\frac12,\qquad x_3=\frac23,
            \qquad 0<\rho<\frac1{64},
\]
and define disjoint even pairs of bumps
\[
 B_j(u)=\rho^{-1}\left\{
       k\left(\frac{u-x_j}{\rho}\right)
          +k\left(\frac{u+x_j}{\rho}\right)\right\}.
\]
Writing $m_i=\int u^i k(u)\dd u$, the column indexed by $j$
of the matrix of moments of orders zero, two and four is
\[
 2\begin{pmatrix}
 1\\ x_j^2+\rho^2m_2\\
 x_j^4+6\rho^2m_2x_j^2+\rho^4m_4
 \end{pmatrix}.
\]
Subtracting suitable multiples of the preceding rows gives its
determinant exactly as
\[
                 8\prod_{i<j}(x_j^2-x_i^2)>0.
\]
Choose the three coefficients of $\sum_j c_jB_j$ to cancel
the corresponding moments of $\psi_0$.  The sum
$\psi=\psi_0+\sum_jc_jB_j$ has the asserted support and plateau;
odd moments vanish by symmetry.  All derivative and integral norms
of this single fixed function are finite.  This is a signed test
function; positivity is neither asserted nor used.

Let $\psi_b(\phi)=\psi((\phi-\widehat\theta)/b)$, interpreted
in the local angular interval containing $\widehat\theta$.
For sufficiently small $\epsilon$, the atom at $\theta$ lies
in the plateau because $C_0\epsilon<b/8$, while the opposite
atom lies outside its support.  Distributional integration by
parts gives
\[
       -\int H(\psi_b''+\psi_b)
             =\ell-\int r\psi_b.
\]
Taylor's formula for $r\in C^{4,\beta}$ about
$\widehat\theta$, with the five vanishing moments, gives
$|\int r\psi_b|\le C M b^{5+\beta}$.
The replacement of $H$ by $\widehat H$ costs at most
\[
 \epsilon\bigl(b^{-1}\|\psi''\|_1+b\|\psi\|_1\bigr).
\]
At the stated scale both leading errors have order
$\epsilon^{(5+\beta)/(6+\beta)}$.  A numerical integration
error bounded by a fixed multiple of $\epsilon/b$ is of this
same order and may be reserved in advance.
\end{proof}

The angular search in Lemma~\ref{lem:two-field-stable-chord}
supplies the $O(\epsilon)$ location error required here.  Its sine
second difference has nonnegative background at every atom-free
interval, whereas an atom within one sixteenth of the interval
radius contributes at most $-c\ell b_0$.  Choosing
$b_0=L\epsilon$ makes that contribution exceed the
$4\epsilon$ data error and the $O(b_0^2)$ smooth background.
The positive lower bound for the chord's vertical projection
separates the two opposite angular neighborhoods uniformly.

For completeness, the smoothing used after mass extraction is the
periodic convolution $T_b$ with
\[
 K_b=\frac85 k_b-\frac45 k_{2b}
                  +\frac8{35}k_{3b}-\frac1{35}k_{4b},
                 \qquad k_b(u)=b^{-1}k(u/b).
\]
Its integral is one and its nonconstant moments through degree six
vanish.  Taylor's formula gives
\[
 \|T_bf-f\|_{C^j}\le Cb^{s-j}\|f\|_{C^{6,\beta}},
   \qquad
 \|T_bq\|_{C^j}\le Cb^{-j}\|q\|_\infty,
                 \qquad 0\le j\le2.
\]
The chord support has two point masses as curvature, so its smoothed
$C^2$ norm is $O(b^{-1})$ and its angular derivative has smoothed
$C^2$ norm $O(b^{-2})$.  Combining these two bounds with the
mass error $O(\epsilon^{(s-1)/s})$ and angular error
$O(\epsilon)$ gives the $C^2$ error
$O(\epsilon^{(s-2)/s})$ in
Lemma~\ref{lem:two-field-stable-chord}.
All integrations here act on acquired finite support representations
and fixed smooth kernels.  They use no additional observations.
